\documentclass[11pt, a4paper]{article}

% --- UNIVERSAL PREAMBLE BLOCK ---
\usepackage[a4paper, top=2.5cm, bottom=2.5cm, left=2cm, right=2cm]{geometry}
\usepackage{fontspec}

\usepackage[english, bidi=basic, provide=*]{babel}

\babelprovide[import, onchar=ids fonts]{english}

% Set default/Latin font to Serif for an academic/publication-grade style
\babelfont{rm}{Noto Serif}

% --- MATHEMATICS PACKAGES ---
\usepackage{amsmath}
\usepackage{amssymb}
\usepackage{amsthm}

% --- THEOREM ENVIRONMENTS ---
\newtheorem{theorem}{Theorem}[section]
\newtheorem{definition}[theorem]{Definition}
\newtheorem{lemma}[theorem]{Lemma}
\newtheorem{observation}[theorem]{Observation}

% --- TITLE AND AUTHOR ---
\title{The Arithmetic Site and Persistent Homology: \\ Prime Emergence as Dimensional Expansion over $\mathbb{F}_1$}
\author{} % Add your name here
\date{\today}

\begin{document}

\maketitle

\begin{abstract}
This paper presents a novel topological framework for the emergence of prime numbers. Rather than viewing the integers ($\mathbb{Z}$) as an \textit{a priori} one-dimensional continuum, we model $\mathbb{Z}$ as the emergent boundary where two distinct geometric realities intersect: the purely multiplicative lattice of the Arithmetic Site over the field with one element ($\mathbb{F}_1$), and the one-dimensional additive continuum. By constructing an Additive-Multiplicative Simplicial Complex and applying persistent homology, we demonstrate that prime numbers function as dimensional basis generators that resolve topological tensions. The persistent $H_0$ barcode encodes prime gaps, while $H_1$ cycles encode the braided computational complexity of factorization. We further define an Arithmetic Jones Polynomial to encode factorization paths, revealing that the topological cost of factorization is conserved within the $\mathbb{F}_1$ tangent space of its base prime.

\vspace{0.5cm}
\noindent \textbf{Keywords:} Persistent Homology, Arithmetic Site, Field with One Element, Prime Distribution, Knot Theory, Computational Complexity. \\
\noindent \textbf{MSC 2020:} 14G10, 11N05, 55N31.
\end{abstract}

\section{Introduction: The Dual Generators of $\mathbb{Z}$}

In standard arithmetic, the integers form a ring equipped with addition and multiplication. However, this structure obscures the profound asymmetry between the two operations. We propose a geometric re-interpretation based on the structural roles of the operational units:

\begin{enumerate}
    \item \textbf{The Additive Identity ($0$):} The singular origin of the space.
    \item \textbf{The Additive Unit ($1$):} A strict \textit{translational} operator that shifts the space linearly.
    \item \textbf{The Multiplicative Units (Primes $p_i$):} \textit{Rotational} or \textit{dimensional} operators that generate orthogonal lattice axes.
\end{enumerate}

In this framework, the integers are not a line, but a braided topological fabric. The primes serve as the "warp" (structural basis generators forming an infinite-dimensional space), and the additive translations serve as the "weft" (the intertwining thread that collapses these dimensions into a sequence).

\section{The Multiplicative Lattice as an Affine $\mathbb{F}_1$-Space}

We begin by isolating the multiplicative structure of the integers, echoing the approach of Connes and Consani in the Arithmetic Site \cite{Connes2014}.

\begin{definition}[$k$-Smooth Lattices]
Let $p_k$ be the $k$-th prime number. Let $S_k$ denote the set of $k$-smooth numbers (integers whose prime factors are all $\leq p_k$). 
\end{definition}

Multiplicatively, $S_k$ is a free commutative monoid generated by $\{p_1, p_2, \dots, p_k\}$. In the language of absolute algebraic geometry, $S_k$ is exactly the affine space of dimension $k$ over the field with one element, denoted as $\mathbb{A}^k_{\mathbb{F}_1}$.

Over $\mathbb{F}_1$, addition does not exist. The space $S_k$ consists of disconnected points floating in a $k$-dimensional void. There is no natural metric or ordering between $x, y \in S_k$ unless one is a direct multiple of the other.

\section{The Additive-Multiplicative Simplicial Complex ($\mathcal{K}$)}

To transition from the static $\mathbb{F}_1$-geometry to the topology of $\mathbb{Z}$, we introduce the "Additive Oracle"—an exogenous probe that attempts to stitch the $\mathbb{A}^k_{\mathbb{F}_1}$ space into a 1-dimensional continuum. We formalize this interaction as a simplicial complex $\mathcal{K}$.

\subsection{Simplices and Boundary Operators}

\begin{itemize}
    \item \textbf{0-Simplices ($C_0$):} The vertices are the natural numbers, $x \in \mathbb{N}$.
    \item \textbf{1-Simplices ($C_1$):} We define two distinct families of directed edges:
    \begin{itemize}
        \item \textit{Additive Edges ($E_A$):} $\alpha_x$ denotes the directed edge from $x$ to $x+1$.
        \item \textit{Multiplicative Edges ($E_M$):} $\mu_{x, p}$ denotes the directed edge from $x$ to $x \cdot p$, where $p$ is prime.
    \end{itemize}
\end{itemize}

The boundary operator $\partial_1 : C_1 \to C_0$ is defined standardly as target minus source:
\begin{align*}
\partial_1(\alpha_x) &= (x+1) - (x) \\
\partial_1(\mu_{x, p}) &= (x \cdot p) - (x)
\end{align*}

\subsection{The Absence of Distributive 2-Simplices}

Crucially, while we map 2-simplices for the commutativity of addition and the commutativity of multiplication, \textbf{we do not define 2-simplices for the distributive law.} The relation $p(x+1) = px + p$ is treated not as an axiomatic triviality, but as a topological hole (a 1-cycle). This omission is the engine of our topological invariants.

\section{Persistent Homology and the Dimensional Filtration}

We convert the complex $\mathcal{K}$ into a persistence module by defining a filtration function based on the sequential expansion of the $k$-smooth lattices.

\begin{definition}[Filtration Function]
Let the discrete time parameter $\epsilon = k$ represent the index of the primes. Let $\text{ind}(x)$ be the index of the largest prime factor of $x$ (with $\text{ind}(1)=0$). We define the filtration $f: \mathcal{K} \to \mathbb{R}$ as:
\begin{itemize}
    \item $f(x) = \text{ind}(x)$ for vertices.
    \item $f(\mu_{x, p_i}) = \max(f(x), i)$ for multiplicative edges.
    \item $f(\alpha_x) = \max(f(x), f(x+1))$ for additive edges.
\end{itemize}
\end{definition}

This yields a sequence of subcomplexes $\mathcal{K}_0 \subseteq \mathcal{K}_1 \subseteq \dots \subseteq \mathcal{K}$, where $\mathcal{K}_k$ is the geometric realization of the interaction between the $+1$ oracle and the $k$-dimensional $\mathbb{F}_1$-lattice.

\subsection{$H_0$ Persistence and Prime Gaps}

At time $\epsilon = k$, the multiplicative edges populate instantly, but the additive edges lag severely because adjacent integers are rarely smooth with respect to the same small set of primes \cite{Stormer1897}.

\begin{theorem}[Prime Gaps as $H_0$ Death Times]
At filtration step $\mathcal{K}_k$, the space is shattered into disconnected $H_0$ components ("floating islands"). The "death" of an $H_0$ component occurs exactly when the lattice expands sufficiently to allow the Additive Oracle to bridge the gap. Consequently, the maximum lifespan of $H_0$ classes in $\mathcal{K}_k$ encodes the bounds of the prime gaps following $p_k$.
\end{theorem}

\section{The Knot Theory of the Arithmetic Braid}

Because we omitted the distributive 2-simplices, a path that leaves the origin multiplicatively and returns additively forms a closed $H_1$ loop. This allows us to map the evaluation of multiplication to formal knot invariants.

\subsection{The Distributive Cycle}

For any integer $x$ and prime $p$, evaluating the multiplicative edge $\mu_{x,p}$ using only the Additive Oracle requires traversing $x(p-1)$ distinct $\alpha$ edges. We formally define the \textbf{Distributive Cycle} $\mathcal{D}(x,p)$ as:
\[
\mathcal{D}(x, p) = \mu_{x, p} - \sum_{i=0}^{x(p-1)-1} \alpha_{x+i}
\]
The length (or "weight") of the additive relaxation of this cycle is exactly $xp - x = x(p-1)$.

\subsection{The Arithmetic Linking Number $\ell(p,q)$}

In a standard symmetric monoidal category, $p \otimes q \cong q \otimes p$. In our computationally braided complex, swapping the order of multiplication induces a topological twist. 

\begin{definition}[Computational Linking Number]
We define the Arithmetic Linking Number $\ell(p,q)$ between two primes as the absolute difference in the additive weights of their braided evaluations:
\[
\ell(p, q) = | p(q-1) - q(p-1) | = | pq - p - pq + q | = | p - q |
\]
\end{definition}

This invariant proves that multiplication is only computationally symmetric when $p = q$. The metric distance $|p-q|$ directly measures the "tightness" of the knot created by swapping the prime fibers in the evaluation braid.

\subsection{The Path Invariant Theorem}

This local knot theory elegantly generalizes to the entire factorization tree of any composite integer $N$. Let $\Gamma = (p_1, p_2, \dots, p_n)$ be an ordered sequence of prime factors evaluating to $N$. 

\begin{theorem}[The Fundamental Path Invariant]
The total additive weight of evaluating any composite integer $N$ depends entirely and exclusively on the initial prime dimension entered. For any factorization braid $\Gamma$, the total additive topological weight is an invariant given by:
\[
W(\Gamma) = N - p_{base}
\]
where $p_{base}$ is the first prime in the operational sequence.
\end{theorem}

\begin{proof}
Let $\Gamma = (p_1, p_2, \dots, p_n)$ be the factorization braid. The total additive weight $W(\Gamma)$ is the sum of the additive relaxations of each intermediate evaluation step. 
\begin{align*}
W(\Gamma) &= \sum_{i=1}^{n-1} \left( \prod_{j=1}^{i} p_j \right) (p_{i+1} - 1) \\
&= p_1(p_2 - 1) + p_1p_2(p_3 - 1) + \dots + \left(\frac{N}{p_n}\right)(p_n - 1)
\end{align*}
Expanding the terms, we observe a telescopic cancellation:
\begin{align*}
W(\Gamma) &= (p_1p_2 - p_1) + (p_1p_2p_3 - p_1p_2) + \dots + \left(N - \frac{N}{p_n}\right) \\
&= N - p_1
\end{align*}
Since $p_1$ is defined as the base prime dimension $p_{base}$, we have $W(\Gamma) = N - p_{base}$, concluding the proof.
\end{proof}

\subsection{The Arithmetic Jones Polynomial and Skein Relations}

To formally classify these factorization braids, we construct an analogue to the Jones polynomial. Let $t$ be a formal variable where the exponent denotes additive weight.

\begin{definition}[State Polynomial of a Braid]
For a factorization braid $\Gamma = (p_1, p_2, \dots, p_n)$ resolving to $N$, with intermediate nodes $N_i = \prod_{j=1}^i p_j$, we define the state polynomial $\mathcal{J}_\Gamma(t)$ as the sum of the localized node weights:
\[
\mathcal{J}_\Gamma(t) = \sum_{i=1}^n t^{N_i}
\]
\end{definition}

In our arithmetic complex, an \textbf{Arithmetic Reidemeister Move} consists of swapping adjacent primes $p_k, p_{k+1}$ in the evaluation sequence (where $k > 1$, preserving $p_{base}$). Let $\Gamma_+$ and $\Gamma_-$ be two braids identical everywhere except at an intermediate junction evaluating from a base $x$. $\Gamma_+$ evaluates $x \to xp \to xpq$, while $\Gamma_-$ evaluates $x \to xq \to xqp$. This yields the \textbf{Arithmetic Skein Relation}:
\[
\mathcal{J}_{\Gamma_+}(t) - \mathcal{J}_{\Gamma_-}(t) = t^{xp} - t^{xq}
\]
Because all intermediate nodes telescope out in the final boundary evaluation, the global knot invariant of the equivalence class of braids collapses to a strictly conserved monomial: $\mathcal{V}_{N, p_{base}}(t) = t^{N - p_{base}}$.

\section{The $\mathbb{F}_1$ Tangent Space and Canonical Roots}

The telescopic cancellation is not merely an algebraic curiosity; it is the manifestation of a conservation law rooted deeply in the differential geometry of the Arithmetic Topos over $\mathbb{F}_1$.

\subsection{The Origin and Tangent Vectors of $\mathbb{F}_1$}

In absolute algebraic geometry, the "origin" of an affine space $\mathbb{A}^k_{\mathbb{F}_1}$ is the multiplicative identity, $1$. To build an integer additively, the Additive Oracle must step away from the origin, requiring the selection of a \textbf{tangent vector}. In our lattice, the prime numbers serve exactly as the basis vectors of the tangent space at the origin $1$, denoted $T_1(\mathbb{A}^\infty_{\mathbb{F}_1})$.

When an evaluation braid $\Gamma$ selects $p_{base}$ as its first operation, it is choosing a canonical tangent vector to enter the multiplicative lattice. Because multiplicative operators commute dimensionally, all subsequent operations are effectively \textbf{isometries} within that specific tangent bundle.

\begin{theorem}[Conservation of Arithmetic Topology]
The topological cost of factorization is conserved within the $\mathbb{F}_1$ tangent space. The distance from the target $N$ to the canonical root $p_{base}$ is invariant under internal permutations because all such paths are homotopic within the tangent bundle $T_{p_{base}}(\mathcal{K})$.
\end{theorem}

\section{The Dimensional Expansion Theorem}

As the Additive Oracle probes the $k$-dimensional $\mathbb{A}^k_{\mathbb{F}_1}$ lattice, it inevitably encounters unresolvable $H_1$ tensions. \textbf{The Additive Oracle acts as the functor of base change $\otimes_{\mathbb{F}_1} \mathbb{Z}$.} To preserve the continuity of the 1-dimensional additive continuum, the space is physically forced to undergo a dimensional expansion. The next prime $p_{k+1}$ is the minimal topological patch—a new orthogonal axis in the $\mathbb{F}_1$ tangent space—injected into the lattice to close the broken $H_1$ cycles and merge the shattered $H_0$ islands.

\section{Conclusion}

By treating the primes as dimensional lattice generators and the $+1$ operation as an exogenous probing oracle, we successfully map the integers into a persistent homology framework. The discovery of the Arithmetic Linking Number, the Arithmetic Skein Relation, and the tangent space conservation law opens immediate avenues for future research, particularly in mapping RSA cryptographic primes to the topological unknotting number of braids across orthogonal $\mathbb{F}_1$ tangent spaces.

\begin{thebibliography}{9}

\bibitem{Connes2014}
Connes, A., \& Consani, C. (2014). \textit{Geometry of the Arithmetic Site}. Advances in Mathematics, 291, 274-329.

\bibitem{Edelsbrunner2010}
Edelsbrunner, H., \& Harer, J. (2010). \textit{Computational Topology: An Introduction}. American Mathematical Society.

\bibitem{Stormer1897}
Størmer, C. (1897). \textit{Quelques théorèmes sur l'équation de Pell $x^2 - Dy^2 = \pm 1$ et leurs applications}. Skrifter Videnskabs-selskabet (Christiania), I, Math.-Naturv. Kl., 2.

\end{thebibliography}

\end{document}